\section*{Declarations}
\paragraph{Data and code availability.} The two locked panels (NAFO 2J3KL spawning stock biomass and catch, 1983--2015, with the 2024 reassessment series for the recent-window checks; Edwards Aquifer J-17 annual mean elevation, recharge and pumping, 1934--2023), the archived campaign results of the two registered protocols, the frozen protocols, the re-derivation core \texttt{p09\_core.py}, the table and figure generators and the verifier \texttt{paper09\_cod\_certification\_v44\_verification.py} are in the source folder \path{paper09_cod_certification_v44_src/} distributed with this manuscript and in the public repository \path{MIKEAA2020/general-sustainability}. No data were collected for this study.
\paragraph{Funding.} None.
\paragraph{Conflicts of interest.} The author declares no conflicts of interest.
\paragraph{Use of generative tools.} Drafting, code and compilation were assisted by a language-model agent under the author's direction; every quantitative statement is reproduced by the shipped code.

\begin{thebibliography}{99}\footnotesize

\bibitem{Abaee2026obstruction} Abaee, A. (2026). An obstruction calculus for viability under incomplete observation. Figshare preprint, version 2. \url{https://doi.org/10.6084/m9.figshare.33716593.v2}

\bibitem{Abaee2026edwforecast} Abaee, A. (2026). Does a one-pool water-balance model improve forecasts of Edwards Aquifer head? A scored test at J-17. Zenodo. \url{https://doi.org/10.5281/zenodo.22552680}

\bibitem{Abaee2026codforecast} Abaee, A. (2026). Does a surplus-production ladder improve forecasts of Northern cod? A scored test on NAFO 2J3KL. Zenodo. \url{https://doi.org/10.5281/zenodo.22553609}

\bibitem{Abaee2026edwintervention} Abaee, A. (2026). Governance operators and viability kernels of the Edwards Aquifer: an intervention-selection test at J-17. Zenodo. \url{https://doi.org/10.5281/zenodo.22553311}

\bibitem{Abaee2026codintervention} Abaee, A. (2026). Robust viability of the 2J3KL limit reference point under a surplus-production map: policy scoring, expansion, and when catch cannot help. Zenodo. \url{https://doi.org/10.5281/zenodo.22552060}

\bibitem{Alley1999} Alley, W.M., Reilly, T.E., Franke, O.L. (1999). \emph{Sustainability of Ground-Water Resources}. U.S. Geological Survey Circular 1186, Denver, CO.

\bibitem{Alley2004} Alley, W.M., Leake, S.A. (2004). The journey from safe yield to sustainability. \emph{Ground Water} 42(1), 12--16.

\bibitem{Aubin1991} Aubin, J.-P. (1991). \emph{Viability Theory}. Birkh\"auser, Boston.

\bibitem{Bene2001} B\'en\'e, C., Doyen, L., Gabay, D. (2001). A viability analysis for a bio-economic model. \emph{Ecological Economics} 36, 385--396.

\bibitem{BertsekasRhodes1971} Bertsekas, D.P., Rhodes, I.B. (1971). On the minimax reachability of target sets and target tubes. \emph{Automatica} 7, 233--247.

\bibitem{Bertsekas1972} Bertsekas, D.P. (1972). Infinite-time reachability of state-space regions by using feedback control. \emph{IEEE Transactions on Automatic Control} 17(5), 604--613.

\bibitem{Blanchini1999} Blanchini, F. (1999). Set invariance in control. \emph{Automatica} 35, 1747--1767.

\bibitem{Butterworth2007} Butterworth, D.S. (2007). Why a management procedure approach? Some positives and negatives. \emph{ICES Journal of Marine Science} 64, 613--617.

\bibitem{DeLaraDoyen2008} De Lara, M., Doyen, L. (2008). \emph{Sustainable Management of Natural Resources: Mathematical Models and Methods}. Springer, Berlin.

\bibitem{Devlin2005} Devlin, J.F., Sophocleous, M. (2005). The persistence of the water budget myth and its relationship to sustainability. \emph{Hydrogeology Journal} 13, 549--554.

\bibitem{DFO2009} DFO (2009). \emph{A Fishery Decision-Making Framework Incorporating the Precautionary Approach}. Fisheries and Oceans Canada, Ottawa.

\bibitem{DFO2016} DFO (2016). Stock assessment of Northern cod (NAFO Divs.\ 2J3KL) in 2016. \emph{DFO Canadian Science Advisory Secretariat Science Advisory Report} 2016/026.

\bibitem{Doyen2012} Doyen, L., Th\'ebaud, O., B\'en\'e, C., Martinet, V., Gourguet, S., Bertignac, M., Fifas, S., Blanchard, F. (2012). A stochastic viability approach to ecosystem-based fisheries management. \emph{Ecological Economics} 75, 32--42.

\bibitem{DoyenDeLara2010} Doyen, L., De Lara, M. (2010). Stochastic viability and dynamic programming. \emph{Systems \& Control Letters} 59, 629--634.

\bibitem{EAA2023} Edwards Aquifer Authority (2025). \emph{2023 Groundwater Discharge and Usage}, Table 1 (after USGS letter report, 5 April 2024). San Antonio, TX. \url{https://www.edwardsaquifer.org/wp-content/uploads/2025/05/2023-Groundwater-Discharge-and-Usage.pdf}

\bibitem{EAA2024} Edwards Aquifer Authority (2024). \emph{Managing the Edwards Aquifer: Critical Period Management and Index Wells}. San Antonio, TX. \url{https://www.edwardsaquifer.org/groundwater-users/critical-period-drought-management/}

\bibitem{EARIP2021} Edwards Aquifer Recovery Implementation Program (2021). \emph{Habitat Conservation Plan}. Edwards Aquifer Authority, San Antonio, TX.

\bibitem{Hutchings1994} Hutchings, J.A., Myers, R.A. (1994). What can be learned from the collapse of a renewable resource? Atlantic cod, \emph{Gadus morhua}, of Newfoundland and Labrador. \emph{Canadian Journal of Fisheries and Aquatic Sciences} 51, 2126--2146.

\bibitem{Krawczyk2013} Krawczyk, J.B., Pharo, A., Serea, O.S., Sinclair, S. (2013). Computation of viability kernels: a case study of by-catch fisheries. \emph{Computational Management Science} 10, 365--396.

\bibitem{KrawczykPharo2013} Krawczyk, J.B., Pharo, A. (2013). Viability theory: an applied mathematics tool for achieving dynamic systems' sustainability. \emph{Mathematica Applicanda} 41, 97--126.

\bibitem{Lucas1976} Lucas, R.E., Jr. (1976). Econometric policy evaluation: a critique. \emph{Carnegie-Rochester Conference Series on Public Policy} 1, 19--46.

\bibitem{Martinet2007} Martinet, V., Th\'ebaud, O., Doyen, L. (2007). Defining viable recovery paths toward sustainable fisheries. \emph{Ecological Economics} 64, 411--422.

\bibitem{Myers1995} Myers, R.A., Cadigan, N.G. (1995). Was an increase in natural mortality responsible for the collapse of northern cod? \emph{Canadian Journal of Fisheries and Aquatic Sciences} 52, 1274--1285.

\bibitem{Myers1997} Myers, R.A., Hutchings, J.A., Barrowman, N.J. (1997). Why do fish stocks collapse? The example of cod in Atlantic Canada. \emph{Ecological Applications} 7, 91--106.

\bibitem{NRC2015} National Research Council (2015). \emph{Review of the Edwards Aquifer Habitat Conservation Plan: Report 1}. The National Academies Press, Washington, DC. \url{https://doi.org/10.17226/21699}

\bibitem{Oubraham2018} Oubraham, A., Zaccour, G. (2018). A survey of applications of viability theory to the sustainable exploitation of renewable resources. \emph{Ecological Economics} 145, 346--367.

\bibitem{Pereau2018} Pereau, J.-C., Mouysset, L., Doyen, L. (2018). Groundwater management in a food security context. \emph{Environmental and Resource Economics} 71(2), 319--336.

\bibitem{Pereau2019} Pereau, J.-C., Pryet, A., Rambonilaza, T. (2019). Optimality versus viability in groundwater management with environmental flows. \emph{Ecological Economics} 161, 109--120.

\bibitem{Puente1978} Puente, C. (1978). \emph{Method of Estimating Natural Recharge to the Edwards Aquifer in the San Antonio Area, Texas}. U.S. Geological Survey Water-Resources Investigations Report 78-10, Austin, TX.

\bibitem{Punt2016} Punt, A.E., Butterworth, D.S., de Moor, C.L., De Oliveira, J.A.A., Haddon, M. (2016). Management strategy evaluation: best practices. \emph{Fish and Fisheries} 17, 303--334.

\bibitem{RegnierDeLara2015} Regnier, E., De Lara, M. (2015). Robust viable analysis of a harvested ecosystem model. \emph{Environmental Modeling \& Assessment} 20, 687--698.

\bibitem{Regular2025} Regular, P.M., et al. (2025). Assessment of the Northern cod stock in NAFO Divisions 2J3KL in 2024. \emph{DFO Canadian Science Advisory Secretariat Research Document} 2025/048.

\bibitem{Rouge2013} Roug\'e, C., Mathias, J.-D., Deffuant, G. (2013). Extending the viability theory framework of resilience to uncertain dynamics, and application to lake eutrophication. \emph{Ecological Indicators} 29, 420--433.

\bibitem{SaintPierre1994} Saint-Pierre, P. (1994). Approximation of the viability kernel. \emph{Applied Mathematics and Optimization} 29, 187--209.

\bibitem{Scanlon2003} Scanlon, B.R., Mace, R.E., Barrett, M.E., Smith, B. (2003). Can we simulate regional groundwater flow in a karst system using equivalent porous media models? Case study, Barton Springs Edwards aquifer, USA. \emph{Journal of Hydrology} 276, 137--158.

\bibitem{Schijns2021} Schijns, R., Froese, R., Hutchings, J.A., Pauly, D. (2021). Five centuries of cod catches in Eastern Canada. \emph{ICES Journal of Marine Science} 78, 2675--2683.

\bibitem{Sophocleous1997} Sophocleous, M. (1997). Managing water resources systems: why ``safe yield'' is not sustainable. \emph{Ground Water} 35(4), 561.

\bibitem{Sophocleous2000} Sophocleous, M. (2000). From safe yield to sustainable development of water resources---the Kansas experience. \emph{Journal of Hydrology} 235, 27--43.

\bibitem{Steinemann2003} Steinemann, A.C. (2003). Drought indicators and triggers: a stochastic approach to evaluation. \emph{Journal of the American Water Resources Association} 39, 1217--1233.

\bibitem{TWDB} Texas Water Development Board. Water Data for Texas, well 6837203 (J-17). \url{https://waterdatafortexas.org/groundwater/well/6837203}

\bibitem{Umphres2025} Umphres, G.D., Choi, N.J. (2025). \emph{Estimated Annual Recharge to the Edwards Aquifer in the San Antonio Area, by Stream Basin or Ungaged Area, 1934--2024}. U.S. Geological Survey data release. \url{https://doi.org/10.5066/P1BI62NY}

\bibitem{Watkins1997} Watkins, D.W., Jr., McKinney, D.C. (1997). Finding robust solutions to water resources problems. \emph{Journal of Water Resources Planning and Management} 123, 49--58.

\end{thebibliography}
\end{document}
